% GUIDE D'UTILISATION - college-eval.sty v2.0
% ==============================================

% ========================================
% 1. STRUCTURE DE BASE D'UN CONTRÔLE
% ========================================

% Fichier principal (élève ou corrigé) :
% --------------------------------------
% !TeX program = lualatex
\documentclass[11pt,a4paper]{article}

% Pour la version élève :
\usepackage{college-eval}

% Pour la version corrigé :
% \usepackage[corrige]{college-eval}

% Configuration du contrôle
\examsetup{6\textsuperscript{e}}{Mathématiques}{Titre du contrôle}{1}{55 minutes}{Calculatrice autorisée}
%          └─ Niveau         └─ Matière       └─ Titre           └─N° └─ Durée    └─ Calculatrice

\begin{document}
	\CorrigeWatermark  % Filigrane "CORRIGÉ" (visible uniquement en mode corrigé)
	\ExamBanner        % Bandeau bleu en haut
	\StudentInfoBox    % Cartouche nom/prénom/classe/date/note
	
	\smallspace
	
	% Contenu du contrôle (fichier séparé recommandé)
	\input{nom_du_fichier_body.tex}
	
\end{document}

% ========================================
% 2. NOUVELLES FONCTIONS DISPONIBLES
% ========================================

% --- Barème automatique ---
\bareme{
	\exbadge{1}{6} \quad      % Exercice 1 : 6 points
	\exbadge{2}{4} \quad      % Exercice 2 : 4 points
	\exbadge{3}{8} \quad      % Exercice 3 : 8 points
	\bonusbadge{2}            % Bonus : +2 points
}

% --- Template pour exercice rapide ---
\ExTemplate{1}{Titre de l'exercice}{6}{book}{
	% Contenu de l'exercice ici
	% Icônes disponibles : book, drafting-compass, calculator, lightbulb, 
	%                      divide, percentage, map-marker-alt, arrows-alt, etc.
}

% --- Auto-évaluation (en fin de contrôle) ---
\AutoEval

% --- Espacement adaptatif ---
\exspace     % Espacement entre exercices (0.3cm)
\smallspace  % Petit espacement (0.2cm)

% ========================================
% 3. BOÎTES DISPONIBLES
% ========================================

% --- Boîte exercice standard (bleue) ---
\begin{exercisebox}{\faIcon{book} Exercice 1 : Titre \points{6}}
	Contenu de l'exercice
\end{exercisebox}

% --- Boîte QCM (bleu clair) ---
\begin{qcmbox}{\faIcon{list-ul} QCM \points{4}}
	Pour chaque question, écris la lettre de la réponse correcte :
	
	\begin{qcmtable}
		\textbf{1.} Question ici ? & Réponse A & Réponse B & Réponse C & \ans{B} \\
		\hline
		\textbf{2.} Autre question ? & Réponse A & Réponse B & Réponse C & \ans{A} \\
		\hline
	\end{qcmtable}
\end{qcmbox}

% --- Boîte construction (verte) ---
\begin{constructionbox}{\faIcon{drafting-compass} Construction \points{7}}
	Construis un triangle ABC tel que...
\end{constructionbox}

% --- Boîte problème (rouge clair) ---
\begin{problembox}{\faIcon{search} Problème \points{8}}
	Un cycliste parcourt...
\end{problembox}

% --- Boîte consignes (grise) ---
\begin{consignebox}
	\faIcon{info-circle} \textbf{CONSIGNES IMPORTANTES}
	\begin{itemize}[leftmargin=*,noitemsep]
		\item[\faIcon{check}] Justifiez vos réponses
		\item[\faIcon{times}] Calculatrice interdite
	\end{itemize}
\end{consignebox}

% ========================================
% 4. GESTION DES RÉPONSES (ÉLÈVE/CORRIGÉ)
% ========================================

% --- Réponse simple (pointillés → texte en rouge) ---
La capitale de la France est \ans{Paris}.

% --- Réponse avec espace blanc personnalisé ---
Réponse : \blanks{3cm}

% --- Zone de réponse multiligne ---
\anslines[0.5cm]{4}{
	% Ce contenu apparaît en ROUGE dans le corrigé
	% En mode élève : 4 lignes de pointillés avec interligne de 0.5cm
	Voici la réponse complète avec plusieurs lignes...
	On peut y mettre des équations :
	\begin{align*}
		x + 5 &= 12 \\
		x &= 7
	\end{align*}
}

% --- Contenu visible uniquement en corrigé ---
\onlycorrige{
	Ce texte n'apparaît que dans le corrigé !
}

% --- Contenu visible uniquement pour l'élève ---
\onlyeleve{
	Ce texte n'apparaît que dans la version élève !
}

% ========================================
% 5. OPÉRATIONS POSÉES (avec xlop)
% ========================================

% Ces commandes affichent l'opération avec :
% - Version élève : résultat en blanc (invisible)
% - Version corrigé : résultat en noir (visible)

\begin{center}
	\AddPose{345}{78}     % Addition posée
\end{center}

\begin{center}
	\SubPose{543}{287}    % Soustraction posée
\end{center}

\begin{center}
	\MulPose{123}{45}     % Multiplication posée
\end{center}

\begin{center}
	\DivEuc{547}{12}      % Division euclidienne posée
\end{center}

% ========================================
% 6. ZONES DE RÉPONSE PERSONNALISÉES
% ========================================

% --- Zone de réponse avec hauteur fixe ---
\zonereponse{4cm}  % Crée une zone grise de 4cm de hauteur

% --- Zone de réponse avec lignes ---
\zonereponseLignes[0.8cm]{5}  % 5 lignes avec interligne de 0.8cm

% ========================================
% 7. ICÔNES FONTAWESOME DISPONIBLES
% ========================================

\faIcon{book}              % Livre
\faIcon{calculator}        % Calculatrice
\faIcon{drafting-compass}  % Compas
\faIcon{lightbulb}         % Ampoule
\faIcon{divide}            % Division
\faIcon{percentage}        % Pourcentage
\faIcon{search}            % Recherche
\faIcon{list-ul}           % Liste
\faIcon{check}             % Coche
\faIcon{times}             % Croix
\faIcon{info-circle}       % Info
\faIcon{clock}             % Horloge
\faIcon{user}              % Utilisateur
\faIcon{calendar}          % Calendrier
\faIcon{star}              % Étoile
\faIcon{arrow-right}       % Flèche droite
\faIcon{map-marker-alt}    % Marqueur
\faIcon{arrows-alt}        % Flèches

% Pour voir plus d'icônes : http://mirrors.ctan.org/fonts/fontawesome5/doc/fontawesome5.pdf

% ========================================
% 8. GRILLE DE COMPÉTENCES (optionnel)
% ========================================

\CompetencesGrid{
	\item[\faIcon{check}] Chercher : Extraire les informations utiles
	\item[\faIcon{check}] Modéliser : Traduire en langage mathématique
	\item[\faIcon{check}] Raisonner : Justifier une démarche
	\item[\faIcon{check}] Calculer : Effectuer des calculs
	\item[\faIcon{check}] Communiquer : Rédiger une solution
}

% ========================================
% 9. EXEMPLE COMPLET D'EXERCICE
% ========================================

\begin{exercisebox}{\faIcon{calculator} Exercice 3 : Calculs \points{8}}

\textbf{1.} Calcule $25 + 37$ : \ans{62} \points{1}

\exspace

\textbf{2.} Effectue la division posée de 156 par 12. \points{3}

\begin{center}
	\DivEuc{156}{12}
\end{center}

\exspace

\textbf{3.} Résous l'équation $2x + 5 = 17$. \points{4}

\anslines[0.5cm]{4}{
	\begin{align*}
		2x + 5 &= 17 \\
		2x &= 17 - 5 \\
		2x &= 12 \\
		x &= 6
	\end{align*}
	La solution est $x = 6$.
}

\end{exercisebox}

% ========================================
% 10. BONNES PRATIQUES
% ========================================

% 1. Créer 3 fichiers pour chaque contrôle :
%    - controle_XX_eleve.tex (version élève)
%    - controle_XX_corrige.tex (version corrigé)  
%    - controle_XX_body.tex (contenu partagé)

% 2. Utiliser \ans{} pour les réponses courtes
%    et \anslines{}{} pour les réponses longues

% 3. Grouper les exercices par type avec les bonnes boîtes :
%    - exercisebox : exercices standard
%    - qcmbox : questionnaires
%    - constructionbox : géométrie
%    - problembox : problèmes complexes

% 4. Toujours mettre \AutoEval à la fin

% 5. Compiler avec LuaLaTeX (support UTF-8 natif + fontawesome5)

% ========================================
% 11. COMPILATION
% ========================================

% Dans le terminal :
% lualatex controle_XX_eleve.tex
% lualatex controle_XX_corrige.tex

% Ou dans votre éditeur LaTeX, configurer :
% !TeX program = lualatex

% ========================================
% FIN DU GUIDE
% ========================================
